\documentclass[12pt, titlepage, a4paper]{article}
\usepackage[a4paper, inner=2.5cm, outer=2.5cm, top=2.25cm, bottom=2.25cm]{geometry} %Lepša postavitev strani
%matematika
\usepackage{amsmath}
\usepackage{amsthm}
\usepackage{amssymb}
\usepackage{empheq}
%
\usepackage[slovene]{babel}
\usepackage{hyperref}
\usepackage{nameref}


% \newtheorem{theorem}{Izrek}
% \newtheorem{lemma}{Lema}
% \newtheorem*{remark}{Opomba}

% \newcommand{\p}{\mathbb{P}}
% \newcommand{\N}{\mathbb{N}}
% \newcommand{\la}[2]{\lambda_{#1} (#2)}
% \newcommand{\pnk}[2]{{#1}^{(#2)}}
% \newcommand{\dnk}{{2}^{(k)}}
% \newcommand{\lak}{\(\lambda_k\)}
% \newcommand{\z}{\zeta}

\title{Implementacija linearnih tipov}
\author{Vida Mlinar}
\date{}


\begin{document}
    
\section{Kaj so linearni tipi}

\subsection{Zakaj tipi}
Računalnik lahko v spomin zapiše objekte v binarnem zapisu. To pomeni, da je vse kar je shranjeno v računalniku, shranjeno z naravnimi števili. To pomeni da je `a' v kodianju utf-8 `01100001' `b' pa je `01100001'. Potem je a + 1 = b. Tega ponavadi nočemo. Za znak + želimo da se obnaša različno glede na to, kakšen argument dobi. 1 + 1 = 2, a + b = ab, a+1 = error. Zato je koristno, da programski jezik ločuje med različnimi tipi objektov. Ta ideja se je razvila že zelo zgodaj v računalništvu
% Teorija tipov je način razmišljanja o matematičnme svetu. Vsak objekt ima svoj natanko določen tip. 

\section{podvečnosti}

\section{Posplošitev linearnih tipov -- linearne ženske}

\section{Affini tipi -- bi tipi?}

\section{Premikanje med linearnimi in afinimi -- trans(formacija)}



\end{document}